\input{../common}

\begin{document}
	%<*content>
	\development{algebra}{equa-diff}{Équations différentielles}

 \summary{}
 
	


 \textbf {Restitution des connaissances}
 
  \begin{exercice} 
  Compléter les phrases ci-dessous.
  \begin{enumerate}
  \item  Les solutions sur $ \mathbb{R} $ de l'équation différentielle  $y^{\prime}- a y = 0$ sont les fonctions  du type $ \cdots $
  \item Une équation différentielle linéaire homogène du second   ordre à coefficients constants est une équation
du type $\cdots$
\item Les solutions sur $ \mathbb{R} $ de l'équation différentielle  $y^{\prime\prime}+\omega^2 y = 0$ sont les fonctions de la forme  $ \cdots $
\item L'équation caractéristique de l'équation différentielle $ay^{\prime\prime}+by^{\prime}+c y = 0$ est $ \cdots $
  \end{enumerate}
\end{exercice}
\begin{exercice} 
Répondre par vrai ou faux.
 \begin{enumerate}
  \item L'équation différentielle $-y^{\prime}= y $ a pour solution générale :\\
  \textbf{a)}\; $ f(x)=K\eexp{x} $\hspace*{3cm} \textbf{b)}\; $ f(x)=\eexp{-x} $\hspace*{3cm} \textbf{c)}\; $ f(x)=K\eexp{-x} $
    \item L'équation différentielle $y^{\prime\prime}= -y $ a pour solution générale :\\
  \textbf{a)}\; $ f(x)=A\eexp{x}+ B\eexp{-x}$\hspace*{1cm} \textbf{b)}\; $ f(x)=A\cos(x)+B\sin(x) $\hspace*{1cm} \textbf{c)}\; $ f(x)=\eexp{x} \paren{A\cos x+B\sin x}$
  \item La fonction  $x\mapsto 5\eexp{-2x} $ est solution de l'équation         différentielle :\\
  \textbf{a)}\; $ y^{\prime}= 2y$\hspace*{1cm} \textbf{b)}\; $ y^{\prime}= -5y $\hspace*{1cm} \textbf{c)}\; $ y^{\prime}= -2y$

  \item La fonction $x\mapsto 2-\eexp{4x}$ est solution de l’équation différentielle :\\
  \textbf{a)}\; $ y^{\prime} +4y =8$\hspace*{2cm} \textbf{b)}\; $ y^{\prime} -4y =4$ \hspace*{2cm} \textbf{a)}\; $ y^{\prime} -4y =8$
      \item L'équation différentielle $y^{\prime\prime}= -\tfrac{1}{x^2} $ a pour solution générale :\\
  \textbf{a)}\; $ f(x)=\ln x+C$\hspace*{1cm} \textbf{b)}\; $ f(x)=\tfrac{1}{x} $\hspace*{1cm} \textbf{c)}\; $ f(x)=-K\ln x$
   \end{enumerate}
\end{exercice}
\bigskip

\textbf{Applications de règles ou de méthodes}
 \begin{exercice} 
Résoudre les équations différentielles suivantes.

\textbf{a)}\; $ y^{\prime}=4y $\hspace*{3cm}
\textbf{b)}\; $ y^{\prime}=\frac{3}{2}y $

\textbf{c)}\; $ y^{\prime}+y=0 $\hspace*{3cm}
\textbf{d)}\; $ 3y^{\prime}-y =0$

\end{exercice}

  \begin{exercice} 
  Résoudre les équations différentielles suivantes.

\textbf{a)}\; $y^{\prime\prime } +6y^{\prime}+9y=0 $\hspace*{3cm}
\textbf{b)}\; $y^{\prime\prime } -4y^{\prime} -5y=0 $

\textbf{c)}\; $ 3y^{\prime\prime }+10 y^{\prime} +3y=0  $\hspace*{3cm}
\textbf{d)}\; $ y^{\prime\prime } +7y^{\prime} =0 $

\end{exercice}
  \begin{exercice} 
  Résoudre  sur $ \mathbb{R} $ les équations différentielles  :\;\;\textbf{a)} \, $y'+4=0$  \qquad \textbf{b)} \, $y"=4$ \qquad \textbf{c)} \, $y"=x$.
\end{exercice}
  \begin{exercice} 
  \begin{enumerate}
\item  Résoudre  dans $ \mathbb{R}$ l'équation différentielle (E) \;: $ 3y'+2y=0 $.
\item Déterminer la solution  g de (E)  dont la courbe représentative  admet pour tangente au point d'abscisse $ 0 $, la droite d'équation \; $ y=-2x+3 $.
\end{enumerate}


\end{exercice}

  \begin{exercice} 
  \begin{enumerate}
\item
Résoudre  dans $ \mathbb{R}$ l'équation différentielle (E) \;: $ -\frac{1}{2}y''+\frac{3}{2}y'-y=0 $
\item Déterminer la solution $ g $ de (E)   dont la courbe représentative passe par le point A$ (0,1) $  et   dont la tangente en ce point    est parallèle à l'axe des abscisses.
\end{enumerate}

\end{exercice}

  \begin{exercice} 
  On considère l'équation différentielle (E) suivante :\;  $y''+2y'+y=x+3$.


 Déterminer les réels $ a$ et $b $  pour  que la fonction $ h $ : $ x\longmapsto ax+b $, soit solution de (E).
 
\end{exercice}
  \begin{exercice} 
  Soit  la fonction $ f $ définie sur $ \mathbb{R} $ par  : $ x\longmapsto (1 + x)e ^{-2x}$.
\begin{enumerate}
\item  Déterminer les  réels $ a $ et $ b $ pour
que $ f $ soit solution de l'équation $ y" + ay' + by = 0$.
\item En déduire  les primitives de $ f $  sur $ \mathbb{R} $.
\end{enumerate}
 
\end{exercice}
  \begin{exercice} 
   Soient les équations différentielles : \;  (E$_{0})$ : \;$  y'+y=0 $\; et \; (E): \;$  y'+y=\eexp{-x}\cos x $.
 \begin{enumerate}
\item  Trouver les réels $ a$ et $b $ pour que la fonction  $ h $ :  $ x\longmapsto\paren{a\cos x+b\sin x}\eexp{-x} $  soit solution de (E).
\item Démontrer  qu'une fonction $f$ est solution de (E) si et seulement si $ f-h$  est solution de (E$_{0})$.
\item Résoudre (E$_{0})$.
\item Déduire des questions précédentes la solution générale de (E).
\item Déterminer la solution   $ g $ de (E)  telle que $ g(0)=1 $.
\end{enumerate}

\end{exercice}
  \begin{exercice} 
  On considère l'équation différentielle (E):\; $ y''+2y'+y=-2\eexp{-x} $. 
\begin{enumerate}
\item Monter  que la fonction $ h$ définie par : $ h(x)=-x^{2}\eexp{-x} $   est une solution particulière de (E).
\item  Résoudre dans $ \mathbb{R}$ l'équation différentielle $ (E') $ \;: $  y''+2y'+y=0 $.
\item Démontrer  que   $ f $  est solution de  (E)   si et seulement si  $ f-h $  est solution de  $ (E') $.
\item En déduire la solution générale  $ f $ de (E).
\item Déterminer la solution $ g $  de (E) satisfaisant aux conditions initiales \;: $g(0)=1 $ et $ g'(0)=0$.
\end{enumerate}

\end{exercice}
  \begin{exercice} 
  On considère sur $ \mathbb{R} $ l'équation différentielle : \;  (E): \;$  y'+2y=\eexp{-2x}$.
 \begin{enumerate}
\item Vérifier que la fonction $ g : x\longmapsto (x+1)\eexp{-2x}$ est une  solution de(E).
\item Démontrer  qu'une   $ f +g$  est solution de (E)   si et seulement si $ f$  est  solution de $y'+2y =0$.
\item Déduire des questions précédentes la solution générale de (E).
\end{enumerate}

\end{exercice}
  \begin{exercice} 
  \begin{enumerate}
\item Résoudre  l'équation différentielle (E): \! $  y''+2y'+5y=0 $.
\item  Déterminer la solution de  $ f $  qui vérifie \; $f(0)=1 $ et $ f'(0)=-1$.
\item On pose\;\; $ F(x)=-\frac{1}{5}\paren{f'(x)+2f(x)} $. 
\begin{enumerate}
\item Démontrer que $ F $  est une primitive de $ f $ sur  $ \mathbb{R} $ puis  expliciter $  F(x) $.
\item En déduire le calcul de $ \int_{0}^{\frac{\pi}{2}}f(x){\rm d}x$
\end{enumerate}
\end{enumerate}

\end{exercice}
  \begin{exercice} 
  Une note de musique est émise en pinçant la corde d'une guitare électrique.
  
La puissance du son émis, initialement de 100 watts, diminue avec le temps t, mesuré en secondes.

On modélise par f(t) la puissance du son émis, exprimée en watt, t secondes après le pincement
de la corde. Le son s'affaiblit à une vitesse proportionnelle à sa puissance, il a été établi
que le coefficient de proportionnalité est de $ - 0,12$.
\begin{enumerate}
\item  Écrire l'équation différentielle traduisant la diminution de son.
\item  Déterminer la fonction f solution de l'équation différentielle (E) qui vérifie la condition
initiale f(0) = 100.
\item Quelle est la puissance du son deux secondes après le pincement de la corde ?
Arrondir au watt près.
\item  Résoudre par le calcul l'équation f(t) = 60, on donnera la valeur exacte
et la valeur approchée à 10$ ^{-3} $ . Interpréter ce résultat
\end{enumerate}
\end{exercice}
  \begin{exercice} 
Un condensateur de capacité $ C $  farads est chargé sous une tension initiale de 20 volts.

 Il se
décharge ensuite dans un résistor de résistance $ R $ ohms.
En notant $ u(t) $ la mesure de la tension en volts au bout de
$ t $ secondes aux bornes du condensateur, $ u $ est alors une
fonction définie sur  $[0 ; +\infty[ $, qui est solution de l'équation  différentielle : \; $ y^{\prime}+\frac{1}{RC}y=0 $
\begin{enumerate}
\item  Résoudre l'équation et en déduire la fonction $ u $.
\item  Pour cette question, $R = 1 000$ et $C = 10^{-4} $.
Pendant combien de temps (au centième de seconde près)
la tension aux bornes du condensateur reste-elle supérieure
ou égale à 5 volts ?
\end{enumerate}
\end{exercice}
  
  
  \bigskip
\textbf {Résolution de problèmes}
 \begin{exercice}
 Le nombre de bactéries $N(t)$ d'une culture initialement à 600 passe au bout de 2 heures à 1 800.
 
On suppose que le taux de croissance est proportionnel au nombre de bactéries présentes.
\begin{enumerate}
\item   Donner une équation différentielle qui traduit le problème puis déterminer $N(t)$ à l'aide des
conditions imposées.
\item  Déterminer le nombre de bactéries après 4 heures.
\item   Déterminer le temps $t$ nécessaire pour que le nombre de bactéries dépasse 12 000.
\end{enumerate}
\end{exercice}
  \begin{exercice} 
  Monsieur Diop est promoteur d'une entreprise agricole. Il a acquis nouvellement une vaste terre plane 
 sur laquelle , il projette y produire de la papaye
Il soumet son projet à un conseil d'ingénieurs pour une étude de marché afin de lui présenter les atouts bénéficiaires sur  ce produit.

Le conseil à la fin de cette étude basée sur un repère orthonormé 
d'unité graphique 1cm pour 100 m, adresse ses solutions à Pierre, un élève compétent en stage auprès du conseil, en ces termes :

<<Le bénéfice à réaliser en milliers de francs en fonction de la quantité $ x $ de papayes en tonnes par an est donné par la fonction $ h $ telle que  $ h(x)^{\prime\prime }-3 h(x)^{\prime} +2h(x)=0  $
et dont la courbe intégrale 
passe par le point  A(0;\; 15000)
et admet en ce point une tangente de coefficient directeur 10000.>>

\medskip
\textbf{Tâche:}\\
  Déterminer le bénéfice maximal annuel à réaliser par l'entreprise de monsieur Diop s'il se lance dans la production de papayes. 
\end{exercice}


  \begin{exercice}
  Le taux de croissance d'une population de girafes dans un parc national peut être modélisé par l'équation différentielle $ \frac{dP}{dt} =0,0002P(200-P)$ où $ t $ est exprimé en années.
  
  
  On sait qu'il y avait 30 girafes quand les scientifiques commençaient à étudier cette population.
  
  
  \textbf{Tâches}\\ Détermine:
  \begin{enumerate}
  \item  Le nombre d'années qu'il faut pour cette population  de girafes atteignent 150.
  \item Quel  maximum  cette population peut-elle atteindre ?
  \end{enumerate}
  \end{exercice}
\end{document}








